\documentclass[10pt,twocolumn]{article}
\usepackage[T1]{fontenc}
\usepackage[utf8]{inputenc}
\usepackage{lmodern}
\usepackage[margin=1.8cm,top=2.4cm,bottom=2cm,columnsep=0.6cm]{geometry}
\usepackage{fancyhdr}
\usepackage{titlesec}
\usepackage{booktabs}
\usepackage{amsmath,amssymb,amsfonts}
\usepackage{microtype}
\usepackage{xcolor}
\usepackage{mdframed}
\usepackage{parskip}
\usepackage{enumitem}
\usepackage{hyperref}

\definecolor{rulered}{RGB}{180,30,30}
\definecolor{boxgray}{RGB}{245,245,245}
\definecolor{keywordgray}{RGB}{100,100,100}

\pagestyle{fancy}
\fancyhf{}
\fancyhead[L]{\textsc{\small Claude Mag}}
\fancyhead[C]{\small\textit{Machine Learning Research Digest}}
\fancyhead[R]{\small 2026-11-11}
\fancyfoot[C]{\small\thepage}
\renewcommand{\headrulewidth}{0.8pt}

\titleformat{\section}{\normalsize\bfseries\color{rulered}}{}{0pt}{}%
  [\vspace{-4pt}\color{rulered}\rule{\columnwidth}{0.4pt}]
\titlespacing{\section}{0pt}{8pt}{4pt}

\hypersetup{colorlinks=true,urlcolor=rulered,linkcolor=black}

\begin{document}
\setlength{\parindent}{0pt}
\setlength{\parskip}{4pt}

{\small\color{keywordgray}\textbf{Keywords:}
  multimodal diffusion LMs \textbullet{}
  hallucination \textbullet{}
  visual grounding \textbullet{}
  training-free inference}\par\vspace{4pt}

{\LARGE\bfseries\raggedright
  VIGOR: Visual Gain Ordering for Hallucination Mitigation in Multimodal
  Discrete Diffusion Language Models}\par\vspace{4pt}

{\large\itshape\raggedright
  Reordering Unmasking by Counterfactual Visual Evidence Consistently Reduces
  Hallucination Without Any Fine-Tuning}\par\vspace{6pt}

{\small\textbf{NeurIPS 2026 Submission}}\par\vspace{2pt}

{\small\color{keywordgray}arXiv:2611.05066 \textbullet{}
  cs.CV / cs.CL \textbullet{} NeurIPS 2026}
\par\vspace{2pt}\noindent{\color{rulered}\rule{\columnwidth}{0.6pt}}\vspace{6pt}

Multimodal discrete diffusion LMs hallucinate because language priors dominate
unmasking confidence, causing tokens to be committed before visual evidence
can influence them. VIGOR reorders unmasking by visual gain --- the difference
in confidence between normal and visually-ablated inference passes --- ensuring
image-grounded tokens are committed first. Applied to LLaDA-V and Lumina-DiMOO,
it consistently reduces hallucination with no training required.

\begin{mdframed}[backgroundcolor=boxgray,linecolor=rulered,linewidth=1pt,
    innerleftmargin=6pt,innerrightmargin=6pt,
    innertopmargin=6pt,innerbottommargin=6pt]
\textbf{\small AT A GLANCE}\par\vspace{4pt}
\begin{description}[leftmargin=0pt,labelsep=4pt,itemsep=2pt,parsep=0pt,
    font=\small\bfseries]
  \item[Problem:] Hallucination in multimodal dLLMs stems partly from premature
    commitment: positions are unmasked by language-prior confidence before
    visual evidence is fully integrated.
  \item[Why it matters:] Multimodal dLLMs are emerging alternatives to
    autoregressive VLMs; their hallucination problem must be solved for
    practical deployment.
  \item[Method:] At each denoising step, compare confidence under normal and
    visually-ablated passes; unmask in order of visual gain (the difference).
  \item[Result:] Consistent hallucination reduction and improved multimodal
    reasoning on both LLaDA-V and Lumina-DiMOO, across all evaluated benchmarks.
  \item[Evidence:] POPE, MMHal-Bench, NocCaps, and standard VQA benchmarks.
\end{description}
\end{mdframed}

\section{The Big Picture}

Discrete diffusion language models generate text by iterative masked-token
prediction: starting from a fully masked sequence, the model at each step
selects the positions it is most confident about and reveals them. When paired
with a vision encoder, these models become multimodal discrete diffusion LMs
(multimodal dLLMs) capable of generating text descriptions conditioned on
images.

Autoregressive VLMs have a well-studied hallucination problem: they generate
plausible-sounding but factually incorrect descriptions of images by over-relying
on language priors. The same problem affects multimodal dLLMs, but with an
additional structural cause specific to their generation process. In an
autoregressive model, every token is generated after all preceding tokens have
been committed, so the full left-context is always available. In a masked dLLM,
the unmasking order is flexible: positions can be committed in any sequence.
The model currently commits positions in order of raw confidence --- but
high-confidence positions are often those that fit well in the sentence based on
language priors alone, independent of the image.

VIGOR identifies this as the \emph{premature commitment problem}: by committing
language-prior-driven positions first, the model forecloses visual evidence
before it can influence downstream positions. VIGOR's solution is to change the
commitment order so that visually-grounded positions are committed first,
providing a visual anchor for subsequent language-driven completions.

\section{Why Existing Approaches Fall Short}

\textbf{Raw confidence ordering} (the standard strategy in masked dLLMs) uses
the model's output probability for the most likely token at each masked position.
This conflates language-prior confidence and visual confidence, unmasking
high-language-prior positions first regardless of their visual grounding.

\textbf{Contrastive decoding methods} for autoregressive VLMs (e.g., VCD, ICD)
subtract language-only scores from image-conditioned scores at the token level,
but they were designed for the per-token generation step in autoregressive models
and do not transfer cleanly to the unmasking-order decision in dLLMs.

\textbf{Training-based approaches} (RL fine-tuning with grounding rewards,
reinforcement from visual feedback) can reduce hallucination but require
additional training data and compute, and do not apply to deployed frozen models.

VIGOR is the first method to apply contrastive visual reasoning to the
\emph{commitment order decision} in masked dLLMs, enabling training-free
hallucination mitigation specific to the dLLM generation structure.

\section{Core Mechanism}

At denoising step $t$, let $m_t$ be the current masking state, $I$ the image,
and $T$ the text context (prompt and already-revealed tokens). For each masked
position $j \in m_t$, VIGOR computes:
\begin{enumerate}[itemsep=2pt,parsep=0pt]
  \item \textbf{Normal confidence:} $P(j \mid I, T) = \max_v p_\theta(v \mid I, T, m_t, j)$,
    the top-1 token probability at position $j$ in the image-conditioned pass.
  \item \textbf{Ablated confidence:} $P(j \mid \emptyset, T) = \max_v p_\theta(v \mid \emptyset, T, m_t, j)$,
    the same probability with the image replaced by a null embedding.
  \item \textbf{Visual gain:} $\text{VG}(j) = P(j \mid I, T) - P(j \mid \emptyset, T)$.
\end{enumerate}
VIGOR selects the top-$k$ positions by visual gain to unmask at step $t$,
rather than the top-$k$ by raw confidence. The null image embedding is a fixed
zero vector in visual token space.

\section{Technical Formulation}

Let $\hat{y}_j = \arg\max_v p_\theta(v \mid I, T, m_t, j)$ be the most likely
token at position $j$ under the image-conditioned model. The visual gain is:
\[
  \text{VG}(j) = p_\theta(\hat{y}_j \mid I, T, m_t, j)
              - p_\theta(\hat{y}_j \mid \emptyset, T, m_t, j).
\]
Positions with $\text{VG}(j) > 0$ have their most likely token more supported
by the image than by language priors alone. Positions with $\text{VG}(j) \approx 0$
are indifferent to the image; positions with $\text{VG}(j) < 0$ are actively
in conflict with the image (language prior predicts something different from
what the image supports).

The unmasking set at step $t$ is:
\[
  \mathcal{U}_t = \arg\text{top-}k_{j \in m_t}\, \text{VG}(j).
\]
The token values assigned to revealed positions are taken from the
image-conditioned pass: $y_j = \hat{y}_j$ for $j \in \mathcal{U}_t$.
The number of positions revealed per step $k$ follows the same schedule as
the base model's standard denoising procedure.

\section{Learning or Inference Procedure}

VIGOR is a pure inference-time method. Required inputs:
\begin{itemize}[itemsep=2pt,parsep=0pt]
  \item A pretrained multimodal dLLM (LLaDA-V, Lumina-DiMOO, or any masked-
    diffusion VLM with accessible token logits).
  \item A null image embedding (zero vector in visual token space; a learned
    null token is also valid and was tested without significant difference).
  \item Two forward passes per denoising step (normal and ablated).
\end{itemize}
The added cost is one extra forward pass per step. For a standard denoising
procedure with $T_d = 10$ steps, VIGOR adds 10 forward passes over baseline,
roughly doubling inference cost.

\section{What the Guarantee Says}

No formal guarantee is claimed. The empirical claim is: on all evaluated
hallucination benchmarks, VIGOR reduces hallucination rates compared to raw-
confidence unmasking ordering. The claim is that visual gain is a better proxy
for visual grounding than raw confidence, and this better proxy translates to
fewer committed hallucinations.

\section{Experimental Findings}

\begin{itemize}[itemsep=2pt,parsep=0pt]
  \item \textbf{POPE:} VIGOR reduces object-level hallucination rate by
    8--14\% (absolute) on both LLaDA-V and Lumina-DiMOO; largest gains on
    adversarial POPE splits.
  \item \textbf{MMHal-Bench:} Hallucination score improves by 0.4--0.7 points
    (on a 0--4 scale) on both backbones.
  \item \textbf{NocCaps:} CIDEr score improves by 3--5 points; caption
    specificity (image-grounded details) increases measurably.
  \item \textbf{Standard VQA:} Performance is unchanged ($\pm 0.5\%$),
    confirming the reordering does not hurt non-hallucination-sensitive tasks.
  \item \textbf{Inference cost:} Roughly $2\times$ baseline; latency is
    acceptable for batch inference applications.
\end{itemize}

\section{Ablations and Interpretation}

\textbf{Null image type:} Zero vector vs.\ learned null token vs.\ Gaussian noise
all produce similar visual gain scores; zero vector is used for simplicity.

\textbf{Top-$k$ selection vs.\ threshold:} Using $\text{VG}(j) > 0$ as a threshold
(rather than a fixed-$k$ selection) gives slightly lower performance because it
occasionally unmasks too few or too many positions per step; fixed-$k$ matching
the base model's schedule is preferred.

\textbf{Position of VIGOR in the denoising schedule:} Applying VIGOR only in the
first half of denoising steps captures 80\% of the hallucination reduction at
half the extra cost, confirming that early commitment decisions are the most
consequential.

\textbf{Comparison with token-level contrastive decoding:} Applying VCD-style
contrastive decoding to the token values (not the commitment order) in masked
dLLMs gives smaller gains than VIGOR, confirming that the commitment-order
effect is the primary mechanism.

\section*{Reference}

NeurIPS 2026 Submission.
\textbf{VIGOR: Visual Gain Ordering for Hallucination Mitigation in Multimodal
Discrete Diffusion Language Models.}
arXiv:2611.05066. NeurIPS 2026.
\url{https://arxiv.org/abs/2611.05066}

\end{document}
